%
% graph.tex -- Template für standalone TIKZ Bilder
%
% (c) 2019 Prof Dr Andreas Müller, Hochschule Rapperswil
%
\documentclass[tikz,12pt]{standalone}
\usepackage{amsmath}
\usepackage{times}
\usepackage{txfonts}
\usepackage{pgfplots}
\usepackage{csvsimple}
\definecolor{darkgreen}{rgb}{0,0.6,0}
\usetikzlibrary{arrows,intersections,math,calc}
\begin{document}
\def\skala{1}
\begin{tikzpicture}[>=latex,thick,scale=\skala]

\pgfmathparse{sqrt(3)+2}
\xdef\rplus{\pgfmathresult}
\pgfmathparse{-sqrt(3)+2}
\xdef\rminus{\pgfmathresult}
\begin{scope}
\clip (-8.1,0) rectangle (8.3,6);
\foreach \s in {-16,...,8}{
	\draw[color=red,line width=1.4pt] plot[domain=0:0.8,samples=40]
		({\s+8*0.5*\rplus*\x*\x},{8*\x});
}
\foreach \s in {-16,...,8}{
	\draw[color=darkgreen,line width=1.4pt] plot[domain=0:0.8,samples=40]
		({\s+8*0.5*(1/\rplus)*\x*\x},{8*\x});
}
\end{scope}
\draw[->] (-8.1,0) -- (8.4,0) coordinate[label={$x$}];
\draw[->] (0,-0.1) -- (0,6.5) coordinate[label={right:$y$}];
\draw (8,-0.1) -- (8,0.1);
\draw (-8,-0.1) -- (-8,0.1);
\node at (8,0) [below] {$1$\strut};
\node at (0,0) [below] {$0$\strut};
\node at (-8,0) [below] {$-1$\strut};
\draw (-0.1,4) -- (0.1,4);
\node at (0,4) [left] {$\frac12$};

\end{tikzpicture}
\end{document}

